% GENERATED FILE. Edit canonical lessons, then run npm run generate:books.
% canonical-source-hash: fnv1a64:d3946ee333297d04
% canonical-lessons: ZH-C328-shenhua, ZH-C328-chuanshuo, ZH-W328-chuan-pass, ZH-C328-liuchuan, ZH-C328-long, ZH-C328-chuanhua

\chapter{A Legend}
\label{ch:zh-a-legend}
\hlchaptermodality{328}

\begin{chapteropening}
\textbf{By the end of this chapter:} \emph{I can talk about myths and legends, say a story was handed down, name a dragon and carry messages.}

Complete the last lesson of chapter 328: I can talk about myths and legends, say a story was handed down, name a dragon and carry messages.
\end{chapteropening}

\section[shénhuà]{shénhuà — myth}

\label{lesson:ZH-C328-shenhua}



\begin{quote}
Before the new one: say the Mandarin for to go into battle, then the Mandarin for a war between countries.
\end{quote}

\subsection*{You'll want to know: shénhuà}
\emph{shénhuà} — \textquotedblleft{}myth\textquotedblright{}.

\begin{quote}
\emph{wǒ xiǎoshíhou tīng guo yí ge shénhuà} — When I was little I heard a myth.
\end{quote}

Its characters come later; for now it is a word for the ear.

\begin{cousinweb}
\emph{shénhuà} is built from \emph{shén} 'a god' and \zh{话} 'talk', both inherited from Middle Chinese, \emph{zyin} and \emph{hwaejH} (Guangyun, in Baxter's transcription), which glosses the first 'a spirit' and the second 'talk'.
\end{cousinweb}

\begin{grammarlens}[title={What you've built}]
One more word for telling what happened.
\end{grammarlens}

\subsection*{Your turn}
\begin{itemize}
\raggedright
  \item \emph{Say it:} \emph{shénhuà}
  \item \emph{Say it:} \emph{shénhuà}, once more
  \item \emph{Say it:} \emph{wǒ xiǎoshíhou tīng guo yí ge shénhuà} — When I was little I heard a myth.
  \item \emph{Recall:} say the Mandarin for to go into battle, then the Mandarin for a war between countries, then say \emph{shénhuà} again
\end{itemize}

\begin{tcolorbox}[breakable,colback=teal!4,colframe=teal!35!black,title={Before you move on}]
What does shénhuà mean? (\textquotedblleft{}Myth\textquotedblright{}.) Say it once more.
\end{tcolorbox}
\section[\texorpdfstring{\zh{传说} (chuánshuō)}{传说 (chuánshuō)}]{\zh{传说} (chuánshuō) — legend; folk tale}

\label{lesson:ZH-C328-chuanshuo}



\begin{quote}
Before the new one: say the Mandarin for a war between countries, then the Mandarin for a myth about the gods.
\end{quote}

\subsection*{You'll want to know: \zh{传说}}
\textbf{\zh{传说}} — \emph{chuánshuō} — \textquotedblleft{}legend; folk tale\textquotedblright{}.

\begin{quote}
\emph{zhège dìfang yǒu yí ge gǔlǎo de chuánshuō} — This place has an old legend.
\end{quote}

The new character, in pieces: \zh{传} — six strokes: the person piece \zh{亻} on the left; on the right two bars, a stroke that runs down, turns right and sweeps down to the left, and a dot. The next lesson puts it in your hand.

\begin{cousinweb}
\textbf{\zh{传说}} (\emph{chuánshuō}) is built from \zh{传} 'to pass on' and \zh{说} 'to tell', both inherited from Middle Chinese, \emph{drjwen} and \emph{sywet} (Guangyun, in Baxter's transcription), which glosses the first 'to turn over, to pass along' and the second 'to tell'.
\end{cousinweb}

\begin{grammarlens}[title={What you've built}]
One more word for telling what happened.
\end{grammarlens}

\subsection*{Your turn}
\begin{itemize}
\raggedright
  \item \emph{Say it:} \emph{chuánshuō}
  \item \emph{Say it:} \emph{chuánshuō}, once more
  \item \emph{Say it:} \emph{zhège dìfang yǒu yí ge gǔlǎo de chuánshuō} — This place has an old legend.
  \item \emph{Recall:} say the Mandarin for a war between countries, then the Mandarin for a myth about the gods, then say \emph{chuánshuō} again
\end{itemize}

\begin{tcolorbox}[breakable,colback=teal!4,colframe=teal!35!black,title={Before you move on}]
What does \zh{传说} mean? (\textquotedblleft{}Legend; folk tale\textquotedblright{}.) Say it once more.
\end{tcolorbox}
\section[\texorpdfstring{\zh{传} (chuán)}{传 (chuán)}]{\zh{传} — the person piece, then two bars and a turning sweep}

\label{lesson:ZH-W328-chuan-pass}



\begin{quote}
The word first, then the character inside it.

\begin{itemize}
\raggedright
  \item \emph{Recall:} say \emph{chuánshuō} — \textbf{R1}, one lesson back, before it decays
\end{itemize}

\begin{quote}
\zh{代} \zh{说}
\end{quote}

\emph{Read it:} each one aloud, with nothing said first — \emph{dài}, \emph{shuō}
\end{quote}

\subsection*{Script — the person piece, then two bars and a turning sweep}
\hlblockfigure{\detokenize{figures/ZH-W328-chuan-pass-filmstrip.pdf}}{How \zh{传} is written, stroke by stroke}

\begin{quote}
\zh{传}
\end{quote}

\textbf{Two pieces, left before right.}

\textbf{The left piece, \zh{亻}, two strokes,} as in \zh{代}.

\textbf{The right piece, four strokes.} A short bar, then a longer bar under it. Then one stroke that starts above the bars, runs down through them, turns right, and sweeps down to the left. Last, a dot at the foot.

2 + 4 = six.

\subsection*{Your turn}
\begin{itemize}
\raggedright
  \item \emph{Write it:} \zh{传}, all six, with the model in front of you
  \item \emph{Cover:} write it again from memory, and count the pieces, not the strokes
\end{itemize}

\begin{tcolorbox}[breakable,colback=teal!4,colframe=teal!35!black,title={Before you move on}]
How many strokes in \zh{传}? (\textbf{Six}.) What is the left piece? (\textbf{\zh{亻}}.) What comes last? (\textbf{A dot at the foot}.)

Source: \href{https://raw.githubusercontent.com/chanind/hanzi-writer-data/68d10a4b21150cae5e1ebbd223eed289cf32d90c/data/\%E4\%BC\%A0.json}{Hanzi Writer Data \zh{传}, pinned stroke paths and medians}.
\end{tcolorbox}
\section[liúchuán]{liúchuán — to be handed down; to spread}

\label{lesson:ZH-C328-liuchuan}



\begin{quote}
Before the new one: say the Mandarin for a myth about the gods, then the Mandarin for an old folk tale.
\end{quote}

\subsection*{You'll want to know: liúchuán}
\emph{liúchuán} — \textquotedblleft{}to be handed down; to spread\textquotedblright{}.

\begin{quote}
\emph{zhège chuánshuō liúchuán le yìqiān nián} — This legend has been handed down for a thousand years.
\end{quote}

Its characters come later; for now it is a word for the ear.

\begin{cousinweb}
\emph{liúchuán} is built from \emph{liú} 'to flow' and \zh{传} 'to pass on', both inherited from Middle Chinese, \emph{ljuw} and \emph{drjwen} (Guangyun, in Baxter's transcription), which glosses the first, quoting an older dictionary, 'water moving'.
\end{cousinweb}

\begin{grammarlens}[title={What you've built}]
One more word for telling what happened.
\end{grammarlens}

\subsection*{Your turn}
\begin{itemize}
\raggedright
  \item \emph{Say it:} \emph{liúchuán}
  \item \emph{Say it:} \emph{liúchuán}, once more
  \item \emph{Say it:} \emph{zhège chuánshuō liúchuán le yìqiān nián} — This legend has been handed down for a thousand years.
  \item \emph{Recall:} say the Mandarin for a myth about the gods, then the Mandarin for an old folk tale, then say \emph{liúchuán} again
\end{itemize}

\begin{tcolorbox}[breakable,colback=teal!4,colframe=teal!35!black,title={Before you move on}]
What does liúchuán mean? (\textquotedblleft{}To be handed down; to spread\textquotedblright{}.) Say it once more.
\end{tcolorbox}
\section[lóng]{lóng — dragon}

\label{lesson:ZH-C328-long}



\begin{quote}
Before the new one: say the Mandarin for an old folk tale, then the Mandarin for to be handed down over time.
\end{quote}

\subsection*{You'll want to know: lóng}
\emph{lóng} — \textquotedblleft{}dragon\textquotedblright{}.

\begin{quote}
\emph{chuánshuō lǐ lóng zhù zài shuǐ lǐ} — In legends, dragons live in water.
\end{quote}

Its characters come later; for now it is a word for the ear.

\begin{cousinweb}
\emph{lóng} is inherited from Middle Chinese \emph{ljowng} (Guangyun, in Baxter's transcription), which calls it the chief of the scaly creatures and quotes an older book: clouds follow the dragon.
\end{cousinweb}

\begin{grammarlens}[title={What you've built}]
One more word for telling what happened.
\end{grammarlens}

\subsection*{Your turn}
\begin{itemize}
\raggedright
  \item \emph{Say it:} \emph{lóng}
  \item \emph{Say it:} \emph{lóng}, once more
  \item \emph{Say it:} \emph{chuánshuō lǐ lóng zhù zài shuǐ lǐ} — In legends, dragons live in water.
  \item \emph{Recall:} say the Mandarin for an old folk tale, then the Mandarin for to be handed down over time, then say \emph{lóng} again
\end{itemize}

\begin{tcolorbox}[breakable,colback=teal!4,colframe=teal!35!black,title={Before you move on}]
What does lóng mean? (\textquotedblleft{}Dragon\textquotedblright{}.) Say it once more.
\end{tcolorbox}
\section[\texorpdfstring{\zh{传话} (chuánhuà)}{传话 (chuánhuà)}]{\zh{传话} (chuánhuà) — to carry messages between people}

\label{lesson:ZH-C328-chuanhua}



\begin{quote}
Before the new one: say the Mandarin for to be handed down over time, then the Mandarin for a dragon.
\end{quote}

\subsection*{You'll want to know: \zh{传话}}
\textbf{\zh{传话}} — \emph{chuánhuà} — \textquotedblleft{}to carry messages between people\textquotedblright{}.

\begin{quote}
\textbf{\zh{他们两个人不说话他在中间传话}} \emph{tāmen liǎng ge rén bù shuōhuà tā zài zhōngjiān chuánhuà} — Those two aren't speaking, so he carries messages between them.
\end{quote}

\begin{cousinweb}
\textbf{\zh{传话}} (\emph{chuánhuà}) is built from \zh{传} 'to pass on' and \zh{话} 'words', both inherited from Middle Chinese, \emph{drjwen} and \emph{hwaejH} (Guangyun, in Baxter's transcription), which glosses the second 'talk'.
\end{cousinweb}

\begin{grammarlens}[title={What you've built}]
One more word for telling what happened.
\end{grammarlens}

\subsection*{Your turn}
\begin{itemize}
\raggedright
  \item \emph{Say it:} \emph{chuánhuà}
  \item \emph{Say it:} \emph{chuánhuà}, once more
  \item \emph{Say it:} \emph{tāmen liǎng ge rén bù shuōhuà tā zài zhōngjiān chuánhuà} — Those two aren't speaking, so he carries messages between them.
  \item \emph{Recall:} say the Mandarin for to be handed down over time, then the Mandarin for a dragon, then say \emph{chuánhuà} again
\end{itemize}

\begin{tcolorbox}[breakable,colback=teal!4,colframe=teal!35!black,title={Before you move on}]
What does \zh{传话} mean? (\textquotedblleft{}To carry messages between people\textquotedblright{}.) Say it once more.
\end{tcolorbox}
